\chapter{Discussion}
\label{ch:discussion}

\section{What the substitution changed}
\label{sec:what-changed}

Replacing score-based diffusion by Riemannian flow matching, with everything
else held fixed, changed three things that can be stated independently of the
outcome of the longer runs.

It changed \emph{what the network is asked to predict}. A score has units of
inverse state and must be rescaled by noise-level-dependent factors before it
can move a pose; a velocity is already in the units of the state. The entire
scalar apparatus that performed that conversion was deleted rather than
adapted (\S\ref{sec:minimal-change}), which is a genuine simplification of the
model's interface, not merely a reorganisation.

It changed \emph{the cost of sampling}. The conditional paths are straight or
geodesic by construction, and \S\ref{sec:nfe} confirms that accuracy is flat
from $5$ to $200$ integration steps --- a fivefold reduction against the
inherited default at no measurable cost. Whether this is an advantage
\emph{over diffusion} awaits the matching sweep on the other arm, and the claim
is stated in the weaker form until then.

And it changed \emph{what is a free parameter}. Under diffusion the source is
determined by the forward process; under flow matching it is a modelling
choice (\S\ref{sec:source-theory}). This is the degree of freedom that
Chapter~\ref{ch:asymmetry} identifies as consequential and
\S\ref{sec:source-results} exploits.

\section{Verdict on the asymmetry hypothesis}
\label{sec:verdict}

\begin{reserved}{Verdict, written against the final evidence}
This section states which of P1--P4 (\S\ref{sec:predictions}) survived. It is
deliberately left unwritten: writing it now would mean committing to a
conclusion before the test. The four anticipated readings and the shape each
would take:

\textbf{(A) Rotation catches up with compute.} P3 refuted. The mechanism is
then a statement about \emph{convergence rate}, not a structural limit: the
time-invariant target makes the rotational problem harder to optimise early,
but not unlearnable. The efficiency result of \S\ref{sec:nfe} becomes the
headline, and the thesis argues that flow matching is competitive above a
compute threshold that should be quantified.

\textbf{(B) The deficit persists.} P3 supported. The thesis argues that the
geodesic path with a Haar source is structurally deficient for
\emph{conditional} generation on $\SOthree$, and that the diffusion noise
schedule functions as an implicit curriculum that the flow-matching
construction discards. \S\ref{sec:source-results} then becomes the causal test
of that claim.

\textbf{(C) Minimal remains worse, informative source repairs it.} P3 and P4
both supported, and the strongest available form of the argument: the
deficiency is diagnosed \emph{and} located in a specific, replaceable
component, with the repair following from the mathematics rather than from
search. Source design is then the thesis's main methodological claim.

\textbf{(D) Accuracy parity, advantage elsewhere.} The asymmetry was a
transient; the substitution is accuracy-neutral but cheaper to sample and, if
\S\ref{sec:confidence-results} succeeds, better equipped for selection. The
thesis then argues on cost and calibration rather than accuracy, with the
explicit caveat that confidence estimation is not flow-specific.

Reserved: \(\approx 400\) words.
\end{reserved}

\section{Is the comparison genuinely controlled?}
\label{sec:controlled}

The strength of any claim here rests on the substitution being isolated, so
the residual differences must be stated rather than assumed away.

In favour: the backbone source is byte-identical between arms, verified by
hash; featurisation, splits, fragmentation, optimiser, schedule and loss
weights are shared; both arms were trained on the same hardware for the same
wall-clock time with the same batch size; and evaluation used one pipeline for
both.

Against: the two arms execute in separate software environments with different
interpreter versions, so a difference attributable to library behaviour cannot
be excluded outright. The diffusion arm is invoked with its own
rotational-score configuration, which is required for it to function but is
nonetheless a difference. And the reconstructed fragmentation differs between
runs for a substantial minority of complexes, which is why the analysis relies
on fragment-level rotational statistics rather than bond-level geometry.

None of these plausibly produces a threefold difference in rotational
enrichment, but they are recorded so the reader can weigh them.

\section{Conceptual differences independent of accuracy}

Some consequences of the substitution do not depend on which model scores
higher. Training is simulation-free and requires no noise schedule to tune, so
one hyperparameter family disappears. Sampling is a deterministic map from
source sample to pose, which makes the sample set a function of $\{x_0^{(k)}\}$
alone and simplifies reasoning about diversity. The learned object is a
well-defined ODE, so exact likelihood is available in principle through the
instantaneous change-of-variables formula --- with the caveat that the
divergence on $\SOthree$ requires care that this thesis has not verified, and
which is therefore listed as future work rather than claimed. And the source
is explicit, which is what makes \S\ref{sec:source-results} possible at all.

\section{Limitations}
\label{sec:limits}

\textbf{Compute.} All runs here are far below the budget behind SigmaDock's
published results. Absolute accuracies are correspondingly low for both arms,
and no statement about benchmark-level performance is made. The comparison is
internally matched, which is the claim that is actually defended.

\textbf{A single training seed per arm.} Run-to-run variability in training is
not characterised, so a difference of the size seen between the two arms
cannot be formally separated from seed variation in training. Sampling-seed
variability \emph{is} characterised and is reported throughout.

\textbf{Evaluation.} PoseBusters is a single benchmark; RMSD-centric metrics
are known to be imperfect proxies for binding-mode correctness; and the
torsional-fingerprint metric explored as a complement proved essentially
uncorrelated with RMSD and failed to compute for a substantial fraction of
molecules, so it is not reported.

\textbf{Inherited design.} The rigid-fragment approximation, the fragmentation
algorithm, and the single-head Newton--Euler reduction are all inherited. In
particular, the last of these means translation and rotation are two
functionals of one output, and a genuinely separate rotational head might
behave differently --- an intervention deliberately excluded here because it is
not flow-specific and would have confounded the comparison.

\textbf{Scope of the mechanism.} The argument of \S\ref{sec:shrinkage} is a
statement about $L^2$ regression under weak conditioning. It predicts that
small fragments, whose rotational readout is most poorly conditioned, should
rotate worst. That prediction is \emph{not} visible in the data
($\Ang{135.8}$ for small fragments against $\Ang{141.7}$ for large ones in
SigmaFlow). With a rotational channel already at chance there is little
dynamic range in which to detect such a trend, so the sub-prediction is
recorded as unconfirmed rather than as supporting evidence.

\section{Future work}

Three directions follow directly. A separate equivariant rotational head would
test whether the two channels compete through their shared readout; this is
the intervention with the largest expected effect and the least flow
specificity, and that combination should be stated plainly. Exact ODE
likelihood is attractive for ranking, but requires first validating the
divergence computation on the manifold against an analytically known test
field. And source design generalises: the argument of
Chapter~\ref{ch:asymmetry} applies to any conditional generative model on a
homogeneous space whose natural invariant measure is uninformative, which
includes several problems in structural biology beyond docking.
